\chapter{Orbital Fragmen dan Molekul Poliatomik}\label{ch:b2:fragment-orbitals}

Struktur Lewis air menunjukkan dua ikatan O--H yang setara dan dua pasangan
elektron bebas yang setara. Namun ketika air diionisasi oleh cahaya
ultraungu, elektron-elektronnya keluar dengan empat energi yang berbeda,
bukan dua; metana, dengan empat ikatan C--H yang setara, menunjukkan dua.
Elektron dalam molekul tidak duduk di dalam ikatan-ikatan di antara
pasangan-pasangan atom: elektron menempati orbital-orbital yang tersebar
di seluruh molekul. Bab ini membangun orbital-orbital itu dari \omterm{def:b2:fragment-orbitals:fragment}{orbital
fragmen}, menjelaskan mengapa air bengkok dan etena datar, serta mengapa
sebuah karbokation lebih menyukai gugus-gugus alkil di sekitarnya.

\begin{recall}
\Cref{ch:b2:diatomic-mos}: \omterm{def:b2:diatomic-mos:lcao}{LCAO}, tumpang-tindih dan simetri, dua tingkat yang
berinteraksi, \omterm{def:b2:diatomic-mos:bonding}{orbital ikatan} dan antiikatan. Jilid Tahun 1: bentuk VSEPR,
hibridisasi sebagai deskripsi, resonansi, karbokation dan urutan
kestabilannya (interaksi ikatan-ikatan C--H dengan orbital kosong sudah
disinggung di sana).
\end{recall}

\section{Metode fragmen}

\begin{definition}[Fragmen]\label{def:b2:fragment-orbitals:fragment}
Yang disebut \emph{fragmen}\index{fragmen} adalah sekelompok atom sebuah
molekul yang ditinjau tersendiri (sebuah atom, sepasang hidrogen, sebuah
gugus \ce{CH2}), dan yang disebut \emph{orbital fragmen}\index{orbital fragmen}
adalah salah satu orbitalnya. Orbital-orbital molekul dibangun dengan
membiarkan orbital-orbital dua fragmen berinteraksi, dua demi dua, seperti
orbital-orbital atom pada molekul diatomik.
\end{definition}

\begin{definition}[Kombinasi teradaptasi simetri]\label{def:b2:fragment-orbitals:symmetry-adapted}
Yang disebut
\emph{kombinasi teradaptasi simetri}\index{kombinasi teradaptasi simetri}
dari orbital-orbital atom yang setara adalah
kombinasi yang oleh setiap operasi simetri molekul (pencerminan terhadap
bidang cermin, rotasi di sekitar sumbu) dibiarkan tidak berubah atau diubah
menjadi lawannya, atau, untuk himpunan yang terdegenerasi, diubah menjadi
kombinasi-kombinasi dari himpunan yang sama.
\end{definition}

\begin{proposition}[Dua hidrogen]\label{prop:b2:fragment-orbitals:h2-pair}
Untuk dua hidrogen setara dengan orbital 1s $h_1$ dan $h_2$, \omterm{def:b2:fragment-orbitals:symmetry-adapted}{kombinasi
teradaptasi simetrinya} adalah $h_1 + h_2$ dan $h_1 - h_2$: yang pertama tidak
berubah, yang kedua berganti tanda, pada setiap operasi yang menukar kedua
atom itu.
\end{proposition}

\begin{proof}
Operasi yang menukar kedua atom mengubah $h_1$ menjadi $h_2$ dan $h_2$
menjadi $h_1$: $h_1 + h_2$ tidak berubah dan $h_1 - h_2$ menjadi
$h_2 - h_1 = -(h_1 - h_2)$.
\end{proof}

\begin{proposition}[Tiga dan empat hidrogen]\label{prop:b2:fragment-orbitals:h4}
Untuk empat hidrogen di sudut-sudut sebuah tetrahedron, kombinasinya
adalah satu kombinasi yang simetris sepenuhnya $h_1 + h_2 + h_3 + h_4$ dan
sehimpunan tiga kombinasi terdegenerasi yang masing-masing mempunyai satu
bidang simpul, berbentuk seperti orbital p. Untuk tiga hidrogen pada sebuah
segitiga (seperti dalam amonia), satu kombinasi simetris dan sepasang
kombinasi terdegenerasi.
\end{proposition}

\admitted

Bentuk kombinasi-kombinasi yang terdegenerasi, dan namanya (a, e, t),
berasal dari teori grup, yang dikembangkan dalam jilid Tahun 3; di sini
nama-nama itu hanyalah label.

\begin{proposition}[Hanya yang sejenis berinteraksi dengan yang sejenis]\label{prop:b2:fragment-orbitals:interaction-rules}
Sebuah \omterm{def:b2:fragment-orbitals:fragment}{orbital fragmen} hanya berinteraksi dengan orbital-orbital \omterm{def:b2:fragment-orbitals:fragment}{fragmen}
lain yang berperilaku sama pada setiap operasi simetri molekul; kekuatan
interaksinya bertambah seiring tumpang-tindih dan berkurang seiring
selisih energinya.
\end{proposition}

\begin{proof}
Jika dua orbital berperilaku berbeda pada suatu pencerminan, yang satu
simetris dan yang lain antisimetris terhadap bidang itu, dan tumpang-tindih
serta \omterm{def:b2:diatomic-mos:integrals}{integral resonansinya} lenyap (\cref{prop:b2:diatomic-mos:symmetry}).
Kebergantungan pada tumpang-tindih dan selisih energi adalah seperti pada
\cref{thm:b2:diatomic-mos:heteronuclear}.
\end{proof}

\begin{method}[Diagram fragmen]\label{met:b2:fragment-orbitals:fragment-diagram}
\begin{enumerate}
  \item Bagilah molekul menjadi dua \omterm{def:b2:fragment-orbitals:fragment}{fragmen}, biasanya atom pusat dan
        himpunan atom di sekitarnya.
  \item Bangunlah \omterm{def:b2:fragment-orbitals:symmetry-adapted}{kombinasi teradaptasi simetri} orbital-orbital luar.
  \item Golongkan orbital-orbital kedua \omterm{def:b2:fragment-orbitals:fragment}{fragmen} menurut perilakunya pada
        bidang-bidang cermin dan sumbu-sumbu molekul.
  \item Biarkan setiap orbital berinteraksi dengan orbital-orbital yang
        perilakunya sama, lebih kuat jika energinya berdekatan; sisanya
        tetap nonikatan.
  \item Isilah dengan elektron valensi lalu bacalah: tingkat terisi
        tertinggi, sifat ikatan, preferensi bentuk.
\end{enumerate}
\end{method}

\section{Air}

Letakkan air pada bidang $yz$, dengan sumbu $z$ membagi dua sudut H--O--H.
Dua bidang cermin memuat sumbu $z$: bidang molekul $yz$ dan bidang $xz$ yang
tegak lurus padanya. Terhadap bidang $xz$, yang menukar kedua hidrogen,
$h_1 + h_2$ simetris dan $h_1 - h_2$ antisimetris. Pada oksigen, 2s dan
$2\mathrm p_z$ simetris terhadap kedua bidang; $2\mathrm p_y$ (pada bidang
molekul, tegak lurus $z$) antisimetris terhadap $xz$; $2\mathrm p_x$
antisimetris terhadap bidang molekul, sesuatu yang tidak dimiliki oleh
kombinasi hidrogen mana pun.

Jadi $h_1 + h_2$ bercampur dengan 2s dan $2\mathrm p_z$, $h_1 - h_2$ dengan
$2\mathrm p_y$, dan $2\mathrm p_x$ tetap nonikatan. Hasilnya enam orbital,
empat di antaranya diisi oleh delapan elektron valensi.

\begin{omfigure}
\begin{tikzpicture}[font=\small]
  % O fragment (left)
  \pic (os) at (0,-2.4) {omlevel={ud}{0.7}}; \node[left] at (-0.45,-2.4) {2s};
  \pic (op1) at (-0.45,0) {omlevel={ud}{0.4}}; \pic (op2) at (0,0) {omlevel={u}{0.4}}; \pic (op3) at (0.45,0) {omlevel={u}{0.4}};
  \node[left] at (-0.7,0) {2p};
  \node at (0,-3.2) {O};
  % H2 fragment (right)
  \pic (hp) at (6,0.9) {omlevel={u}{0.7}}; \node[right] at (6.4,0.9) {$h_1 + h_2$};
  \pic (hm) at (6,1.8) {omlevel={u}{0.7}}; \node[right] at (6.4,1.8) {$h_1 - h_2$};
  \node[anchor=north west] at (6.4,0.6) {fragmen \ce{H2}};
  % MOs
  \pic (m1) at (3,-2.9) {omlevel={ud}{0.7}};  \node[omsymlabel, below=0.17cm, inner sep=0.5pt] at (m1-c) {$2a_1$};
  \pic (m2) at (3,-1.55) {omlevel={ud}{0.7}};  \node[omsymlabel, below=0.17cm, inner sep=0.5pt] at (m2-c) {$1b_2$};
  \pic (m3) at (3,-0.78) {omlevel={ud}{0.7}};  \node[omsymlabel, below=0.17cm, inner sep=0.5pt] at (m3-c) {$3a_1$};
  \pic (m4) at (3,0) {omlevel={ud}{0.7}};  \node[omsymlabel, below=0.17cm, inner sep=0.5pt] at (m4-c) {$1b_1$};
  \pic (m5) at (3,2.4) {omlevel={empty}{0.7}}; \node[omsymlabel, below=0.17cm, inner sep=0.5pt] at (m5-c) {$4a_1$};
  \pic (m6) at (3,3.1) {omlevel={empty}{0.7}}; \node[omsymlabel, below=0.17cm, inner sep=0.5pt] at (m6-c) {$2b_2$};
  \draw[omcorr, densely dashed] (os-r) -- (m1-l) (op3-r) -- (m2-l) (op3-r) -- (m3-l) (op3-r) -- (m4-l) (op3-r) -- (m5-l) (op3-r) -- (m6-l);
  \draw[omcorr, densely dashed] (hp-l) -- (m1-r) (hm-l) -- (m2-r) (hp-l) -- (m3-r) (hp-l) -- (m5-r) (hm-l) -- (m6-r);
\end{tikzpicture}

{\small Diagram \omterm{def:b2:fragment-orbitals:fragment}{fragmen} air, skematis. Kombinasi simetris hidrogen-hidrogen
bercampur dengan orbital 2s dan $2\mathrm p_z$ oksigen (orbital $a_1$),
kombinasi antisimetris dengan $2\mathrm p_y$ ($b_2$); $2\mathrm p_x$, yang
tegak lurus bidang molekul, tetap nonikatan ($1b_1$, pada tingkat 2p
oksigen). Empat tingkat valensi terisi, empat pita fotoelektron; yang
tertinggi, $1b_1$, terionisasi pada \qty{12.62}{eV}. Label
$a_1$, $b_1$, $b_2$ adalah nama (dari teori grup, dalam jilid Tahun 3).}
% ledger: ie:H2O
\end{omfigure}

\begin{proposition}[Mengapa air bengkok]\label{prop:b2:fragment-orbitals:bent-water}
Dalam H--O--H yang linear, dua orbital p oksigen yang tegak lurus sumbu
bersifat nonikatan. Pembengkokan memungkinkan salah satunya, yang terletak
pada bidang molekul, bertumpang-tindih dengan $h_1 + h_2$ dan bercampur
dengan orbital 2s: tingkat terisi itu turun, lebih jauh daripada naiknya
tingkat-tingkat terisi yang lain. Dengan delapan elektron valensi, molekul
yang bengkok lebih stabil.
\end{proposition}

\begin{proof}[Argumen]
Dalam molekul linear ($z$ sepanjang H--O--H), $h_1 + h_2$ mempunyai
tumpang-tindih nol dengan $2\mathrm p_x$ dan $2\mathrm p_y$ (masing-masing
antisimetris terhadap bidang yang memuat sumbu, sedangkan kombinasinya
simetris); keduanya merupakan dua pasangan nonikatan yang terdegenerasi.
Ketika dibengkokkan pada bidang $yz$, $2\mathrm p_y$ sebagian mengarah ke
hidrogen-hidrogen: tumpang-tindihnya dengan $h_1 + h_2$ bertambah dari
nol, dan tingkat terisi yang dibawanya distabilkan (tingkat $3a_1$ molekul
bengkok). Tingkat ikatan $1b_2$ kehilangan sedikit tumpang-tindih dan naik
sedikit. Hasil bersihnya adalah penurunan, yang paling besar jika tingkat
yang turun itu terisi, seperti halnya dengan delapan elektron. Molekul
dengan hanya empat elektron valensi, seperti \ce{BeH2}, membiarkan tingkat
itu kosong dan linear.
\end{proof}

\section{Metana, amonia, dan spektrum fotoelektron}

Dalam metana, keempat orbital hidrogen menghasilkan kombinasi simetris,
yang cocok dengan orbital 2s karbon, dan sehimpunan tiga kombinasi
terdegenerasi, yang cocok dengan ketiga orbital 2p-nya. Empat \omterm{def:b2:diatomic-mos:bonding}{orbital
ikatan} terisi: satu yang rendah ($a_1$) dan tiga yang terdegenerasi pada
energi yang lebih tinggi ($t_2$). Amonia, dengan tiga hidrogen,
mempertahankan sebuah orbital terisi pada nitrogen yang mengarah menjauhi
hidrogen-hidrogen, yaitu pasangan elektron bebasnya, sebagai tingkat
terisi tertinggi.

\begin{omfigure}
\begin{tikzpicture}[font=\small]
  \pic (cs) at (0,-1.6) {omlevel={ud}{0.7}}; \node[left] at (-0.45,-1.6) {2s};
  \pic (cp1) at (-0.45,0) {omlevel={u}{0.4}}; \pic (cp2) at (0,0) {omlevel={u}{0.4}}; \pic (cp3) at (0.45,0) {omlevel={empty}{0.4}};
  \node[left] at (-0.7,0) {2p};
  \node at (0,-2.5) {C};
  \pic (ha) at (6,-0.8) {omlevel={ud}{0.7}}; \node[right] at (6.4,-0.8) {$a_1$};
  \pic (ht1) at (5.55,0.4) {omlevel={u}{0.4}}; \pic (ht2) at (6,0.4) {omlevel={u}{0.4}}; \pic (ht3) at (6.45,0.4) {omlevel={empty}{0.4}};
  \node[right] at (6.7,0.4) {$t_2$};
  \node at (6,-1.6) {fragmen \ce{H4}};
  \pic (m1) at (3,-2.4) {omlevel={ud}{0.7}}; \node[omsymlabel, below=0.17cm, inner sep=0.5pt] at (m1-c) {$1a_1$};
  \pic (m2a) at (2.55,-0.9) {omlevel={ud}{0.4}}; \pic (m2b) at (3,-0.9) {omlevel={ud}{0.4}}; \pic (m2c) at (3.45,-0.9) {omlevel={ud}{0.4}};
  \node[omsymlabel, below=0.17cm, inner sep=0.5pt] at (m2b-c) {$1t_2$};
  \pic (m3) at (3,1.8) {omlevel={empty}{0.7}}; \node[omsymlabel, below=0.17cm, inner sep=0.5pt] at (m3-c) {$2a_1$};
  \pic (m4a) at (2.55,2.5) {omlevel={empty}{0.4}}; \pic (m4b) at (3,2.5) {omlevel={empty}{0.4}}; \pic (m4c) at (3.45,2.5) {omlevel={empty}{0.4}};
  \node[omsymlabel, below=0.17cm, inner sep=0.5pt] at (m4b-c) {$2t_2$};
  \draw[omcorr, densely dashed] (cs-r) -- (m1-l) (cp3-r) -- (m2a-l) (cs-r) -- (m3-l) (cp3-r) -- (m4a-l);
  \draw[omcorr, densely dashed] (ha-l) -- (m1-r) (ht1-l) -- (m2c-r) (ha-l) -- (m3-r) (ht1-l) -- (m4c-r);
\end{tikzpicture}

{\small Diagram \omterm{def:b2:fragment-orbitals:fragment}{fragmen} metana, skematis. Kombinasi simetris keempat
hidrogen cocok dengan orbital 2s karbon, ketiga kombinasi yang
terdegenerasi cocok dengan orbital-orbital 2p-nya: dua tingkat terisi, dan
karena itu dua pita fotoelektron, dengan yang di atas (yang terionisasi
pertama, pada \qty{12.61}{eV}) memuat enam elektron.}
% ledger: ie:CH4
\end{omfigure}

\begin{definition}[Spektroskopi fotoelektron]\label{def:b2:fragment-orbitals:photoelectron}
Yang disebut \emph{spektroskopi fotoelektron}\index{spektroskopi fotoelektron}
mengukur energi kinetik elektron-elektron yang dilepaskan dari molekul
oleh radiasi monokromatik yang energinya diketahui; setiap pita spektrum
bersesuaian dengan pengambilan sebuah elektron dari satu tingkat terisi,
dan posisinya memberikan energi ionisasi tingkat itu.
\end{definition}

\begin{proposition}[Pendekatan Koopmans]\label{prop:b2:fragment-orbitals:koopmans}
Energi yang diperlukan untuk mengambil sebuah elektron dari orbital terisi
kira-kira sama dengan minus \omterm{def:b2:atomic-orbitals:orbital-energy}{energi orbital} itu, $I \approx -\varepsilon$.
\end{proposition}

\admitted

Pendekatan ini mengabaikan relaksasi elektron-elektron lain setelah
ionisasi; pendekatan itu dibenarkan dalam jilid Tahun 3. Pendekatan ini
mengubah spektrum fotoelektron menjadi diagram orbital: air menunjukkan
empat pita valensi, metana dua, sesuai dengan diagram-diagram \omterm{def:b2:fragment-orbitals:fragment}{fragmen} ---
dan bertentangan dengan gambaran dua pasangan elektron bebas yang setara
dan dua ikatan yang setara dalam air, yang memerikan kerapatan elektron
total yang sama dengan cara lain tetapi tidak memberikan energi
elektron-elektron yang diambil.

\section{Etena}

Etena dapat dibangun dari dua \omterm{def:b2:fragment-orbitals:fragment}{fragmen} \ce{CH2}. Setiap \ce{CH2} mempunyai,
selain orbital-orbital kedua ikatan C--H-nya, sebuah orbital yang mengarah
menjauhi hidrogen-hidrogen di dalam bidang (yang membentuk ikatan $\sigma$
C--C dengan pasangannya) dan sebuah orbital p yang tegak lurus bidang.
Kedua orbital p bergabung menjadi $\pi$ yang terisi dan $\pi^*$ yang
kosong.

\begin{omfigure}
\begin{tikzpicture}[font=\small]
  \begin{scope}
    \pic at (0,0) {omp={0.85}{90}}; \pic at (1.2,0) {omp={0.85}{90}};
    \draw[thick] (0,0) -- (1.2,0);
    \draw[thick] (0,0) -- (-0.5,-0.45) (0,0) -- (-0.55,0.35) (1.2,0) -- (1.7,-0.45) (1.2,0) -- (1.75,0.35);
    \node at (0.6,-1.1) {$\pi$ (terisi)};
  \end{scope}
  \begin{scope}[xshift=3.6cm]
    \pic at (0,0) {omp={0.85}{90}}; \pic at (1.2,0) {omp={-0.85}{90}};
    \draw[thick] (0,0) -- (1.2,0);
    \draw[thick] (0,0) -- (-0.5,-0.45) (0,0) -- (-0.55,0.35) (1.2,0) -- (1.7,-0.45) (1.2,0) -- (1.75,0.35);
    \node at (0.6,-1.1) {$\pi^*$ (kosong)};
  \end{scope}
  \begin{scope}[xshift=7.4cm]
    \pic at (0,0) {omp={0.85}{90}}; \pic at (1.2,0) {omp={0.85}{0}};
    \draw[thick] (0,0) -- (1.2,0);
    \node at (0.6,-1.1) {terpuntir $90^\circ$};
    \node[font=\scriptsize, align=center] at (0.6,1.3) {tumpang-tindih nol};
  \end{scope}
\end{tikzpicture}

{\small Sistem $\pi$ etena, skematis (molekul dilihat dari sisi, dengan
hidrogen-hidrogen di dalam bidang). Kedua orbital p yang tegak lurus
bidang bergabung sefase ($\pi$, terisi) dan berlawanan fase ($\pi^*$,
kosong). Memuntir satu \ce{CH2} sebesar $90^\circ$ membuat kedua orbital p
saling tegak lurus: tumpang-tindihnya, dan ikatan $\pi$, lenyap.}
\end{omfigure}

Ikatan $\pi$ menjaga etena tetap datar: memuntir satu ujung memalingkan
orbital-orbital p satu sama lain, tumpang-tindihnya turun sebanding dengan
kosinus sudut puntiran, dan pada $90^\circ$ ikatan $\pi$ hilang. Energi
sebuah ikatan $\pi$ harus dipasok untuk memuntir molekul itu; itulah
sebabnya isomer cis dan trans alkena tidak saling berubah pada suhu kamar,
sedangkan rotasi di sekitar ikatan tunggal bebas.

\section{Hiperkonjugasi}

\begin{definition}[Hiperkonjugasi]\label{def:b2:fragment-orbitals:hyperconjugation}
Yang disebut \emph{hiperkonjugasi}\index{hiperkonjugasi} adalah interaksi
orbital $\sigma$ terisi dari ikatan C--H atau C--C dengan orbital p atau
$\pi^*$ di dekatnya yang kosong atau terisi sebagian dan sejajar dengannya.
\end{definition}

\begin{proposition}[Gugus alkil menstabilkan karbokation]\label{prop:b2:fragment-orbitals:hyperconjugation}
Sebuah karbokation distabilkan oleh setiap ikatan C--H atau C--C pada
karbon tetangga yang dapat sejajar dengan orbital p kosongnya: makin
banyak ikatan seperti itu, makin stabil kationnya; dari situlah urutan
tersier > sekunder > primer > metil.
\end{proposition}

\begin{proof}[Argumen]
Orbital p kosong pada karbon kationik dan orbital $\sigma_{\text{C--H}}$ terisi
di sebelahnya membentuk sistem dua tingkat dengan selisih $\Delta$ yang
besar dan kopling $\beta$ yang kecil (tak nol hanya jika ikatan itu kira-kira
sejajar dengan orbital p). Menurut \cref{thm:b2:diatomic-mos:heteronuclear},
tingkat $\sigma$ yang terisi turun sekitar $\beta^2/\Delta$, dan kedua
elektron memperoleh dua kali nilai itu; tingkat yang kosong naik, tanpa
biaya. Setiap ikatan yang tersejajarkan dengan baik menambahkan
stabilisasinya sendiri; gugus metil selalu mempunyai satu ikatan C--H yang
hampir sejajar, bagaimanapun rotasinya.
\end{proof}

\begin{omfigure}
\begin{tikzpicture}[font=\small]
  \pic at (0,0) {omp={0.9}{90}};
  \node[circle, fill=black, inner sep=1.2pt] at (0,0) {};
  \node[font=\scriptsize, anchor=north west] at (0.08,-0.05) {C$^+$};
  \draw[thick] (0,0) -- (1.4,0);
  \node[circle, fill=black, inner sep=1.2pt] at (1.4,0) {};
  \draw[thick] (1.4,0) -- (1.75,0.85);
  \pic at (1.58,0.42) {omlobe={0.75}{68}};
  \pic at (1.58,0.42) {omlobe={0.45}{248}};
  \node[font=\scriptsize] at (1.95,1.0) {H};
  \draw[thick] (1.4,0) -- (2.0,-0.45) (1.4,0) -- (1.1,-0.6);
  \draw[<->, omProp, thick] (0.35,0.8) .. controls (0.8,1.2) and (1.2,1.2) .. (1.5,0.95);
  \node[font=\scriptsize, align=left, anchor=west] at (2.4,0.4) {cuping $\sigma_{\text{C--H}}$ terisi sejajar\\ dengan orbital p kosong};
  % level scheme
  \begin{scope}[xshift=7.2cm, yshift=-0.4cm]
    \pic (p) at (0,1.4) {omlevel={empty}{0.7}}; \node[left] at (-0.4,1.4) {p kosong};
    \pic (s) at (2.2,-0.8) {omlevel={ud}{0.7}}; \node[right] at (2.6,-0.8) {$\sigma_{\text{C--H}}$};
    \pic (lo) at (1.1,-1.2) {omlevel={ud}{0.7}};
    \pic (hi) at (1.1,1.7) {omlevel={empty}{0.7}};
    \draw[omcorr, densely dashed] (p-r) -- (lo-l) (p-r) -- (hi-l) (s-l) -- (lo-r) (s-l) -- (hi-r);
    \node[font=\scriptsize, anchor=north] at (1.1,-1.35) {distabilkan sebesar $\approx \beta^2/\Delta$};
  \end{scope}
\end{tikzpicture}

{\small \omterm{def:b2:fragment-orbitals:hyperconjugation}{Hiperkonjugasi} dalam karbokation, skematis. \omterm{def:b2:diatomic-mos:bonding}{Orbital ikatan} C--H
terisi pada karbon tetangga, yang sejajar dengan orbital p kosong,
menyerahkan sebagian kerapatan elektronnya kepada orbital itu; tingkat yang
terisi distabilkan seperti pada setiap interaksi dua tingkat dengan energi
yang berbeda.}
\end{omfigure}

\section{Latihan}

\begin{exercise}[$\star$]\label{exo:b2:fragment-orbitals:1}
Tuliskan kedua \omterm{def:b2:fragment-orbitals:symmetry-adapted}{kombinasi teradaptasi simetri} orbital-orbital 1s hidrogen
pada air, lalu nyatakan mana di antaranya yang berganti tanda pada bidang
yang membagi dua sudut H--O--H dan tegak lurus molekul.
\end{exercise}

\begin{exercise}[$\star$]\label{exo:b2:fragment-orbitals:2}
Berapa \omterm{def:b2:diatomic-mos:lcao}{orbital molekul} valensi yang dimiliki air, berapa elektron
valensinya, dan berapa orbital valensi yang terisi?
\end{exercise}

\begin{exercise}[$\star$]\label{exo:b2:fragment-orbitals:3}
Mengapa spektrum fotoelektron metana menunjukkan dua pita valensi, yang
satu tiga kali lebih kuat daripada yang lain?
\end{exercise}

\begin{exercise}[$\star$]\label{exo:b2:fragment-orbitals:4}
Mengapa keenam atom etena terletak pada satu bidang?
\end{exercise}

\begin{exercise}[$\star\star$]\label{exo:b2:fragment-orbitals:5}
Energi ionisasi pertama amonia adalah \qty{10.07}{eV}, energi ionisasi
pertama metana \qty{12.61}{eV}. Orbital mana yang kehilangan elektron dalam
setiap kasus, dan mengapa amonia lebih mudah diionisasi?
\end{exercise}

\begin{exercise}[$\star\star$]\label{exo:b2:fragment-orbitals:6}
Bandingkan tingkat-tingkat terisi air yang linear dan yang bengkok, lalu
jelaskan mengapa molekul seperti \ce{BeH2}, dengan empat elektron valensi,
linear.
\end{exercise}

\begin{exercise}[$\star\star$]\label{exo:b2:fragment-orbitals:7}
Jelaskan, dengan hasil dua tingkat, urutan kestabilan karbokation
\ce{CH3+}, \ce{CH3CH2+}, \ce{(CH3)2CH+}, \ce{(CH3)3C+}.
\end{exercise}

\begin{exercise}[$\star\star$]\label{exo:b2:fragment-orbitals:8}
Tumpang-tindih dua orbital p yang terpuntir sebesar sudut $\theta$ di
sekitar ikatan sebanding dengan $\cos\theta$. Jika \omterm{def:b2:diatomic-mos:integrals}{integral resonansi}
sebanding dengan tumpang-tindih, bagaimana stabilisasi $\pi$ bergantung
pada $\theta$? Pada sudut berapa stabilisasi itu turun menjadi separuhnya?
\end{exercise}

\begin{exercise}[$\star\star$]\label{exo:b2:fragment-orbitals:9}
Boran \ce{BH3} datar, amonia piramidal. Dengan \omterm{def:b2:fragment-orbitals:fragment}{orbital fragmen} sebuah atom
pusat dan tiga hidrogen, jelaskan perbedaannya melalui jumlah elektron
yang dimiliki masing-masing.
\end{exercise}

\begin{exercise}[$\star\star\star$]\label{exo:b2:fragment-orbitals:10}
Ion \ce{H3+} berbentuk segitiga sama sisi. Dengan $\alpha$ dan $\beta$ untuk
orbital-orbital 1s dan tumpang-tindih diabaikan, tentukan ketiga \omterm{def:b2:atomic-orbitals:orbital-energy}{energi
orbitalnya} (kombinasi simetris berenergi $\alpha + 2\beta$), lalu jelaskan
mengapa dua elektron sudah cukup untuk mengikatnya.
\end{exercise}

\begin{exercise}[$\star\star\star$]\label{exo:b2:fragment-orbitals:11}
Sistem alil \ce{CH2=CH-CH2} mempunyai tiga orbital p yang sejajar. Tebaklah
bentuk ketiga orbital $\pi$-nya (tanda pada setiap karbon) dari jumlah
simpulnya, dan mana yang merupakan tingkat terisi tertinggi anion alil.
\end{exercise}

\begin{exercise}[$\star\star\star$]\label{exo:b2:fragment-orbitals:12}
Bangunlah sistem $\pi$ karbon dioksida dari orbital-orbital 2p yang tegak
lurus sumbu O--C--O: berapa orbital $\pi$, mana yang bersifat ikatan,
nonikatan, antiikatan, dan berapa elektron yang dimuatnya?
\end{exercise}

\clearpage
\enlargethispage{3\baselineskip}
\section{Soal: Bentuk Molekul Air}

\begin{problem}[{Soal akhir pekan --- \omterm{def:b2:fragment-orbitals:fragment}{fragmen} oksigen dan \ce{H2},
orbital-orbital H--O--H yang linear, tingkat-tingkat yang bergeser ketika
molekul itu membengkok, dan spektrum fotoelektron dengan pendekatan
Koopmans}]
\label{pb:b2:fragment-orbitals:1}
Data: geometri air, $r(\ce{O-H}) = \qty{95.8}{pm}$, sudut $104.48^\circ$;
energi ionisasi pertama: \ce{H2O} \qty{12.62}{eV}, \ce{O}
\qty{13.62}{eV}. Air terletak pada bidang $yz$, dengan $z$ sepanjang garis
bagi sudut.
% ledger: geo:H2O.rOH, geo:H2O.aHOH, ie:H2O, ie:O

\medskip
\noindent\textbf{Bagian I --- \omterm{def:b2:fragment-orbitals:fragment}{Fragmen}.}
\begin{enumerate}
  \item Berapa elektron valensi yang dimiliki air, dan dari atom-atom mana?
  \item Tuliskan kedua \omterm{def:b2:fragment-orbitals:symmetry-adapted}{kombinasi teradaptasi simetri} orbital-orbital 1s
        hidrogen.
  \item Orbital oksigen manakah yang simetris terhadap kedua bidang cermin?
  \item Orbital oksigen manakah yang antisimetris hanya terhadap bidang
        $xz$?
  \item Orbital oksigen manakah yang tidak mempunyai pasangan di antara
        kombinasi-kombinasi hidrogen?
  \item Berapa \omterm{def:b2:diatomic-mos:lcao}{orbital molekul} yang dihasilkan, dan berapa yang terisi?
\end{enumerate}

\medskip
\noindent\textbf{Bagian II --- H--O--H yang linear.}
\begin{enumerate}[resume]
  \item Dalam molekul linear sepanjang $z$, kombinasi manakah yang
        berinteraksi dengan $2\mathrm p_z$ oksigen?
  \item Kombinasi manakah yang berinteraksi dengan 2s oksigen?
  \item Orbital oksigen manakah yang tetap nonikatan, dan dengan degenerasi
        berapa?
  \item Isilah tingkat-tingkatnya dengan delapan elektron valensi.
  \item Apakah air yang linear akan bersifat \omterm{def:b2:diatomic-mos:magnetism}{paramagnetik}?
  \item Di mana tingkat terisi tertingginya akan terletak, dibandingkan
        dengan orbital 2p oksigen?
\end{enumerate}

\medskip
\noindent\textbf{Bagian III --- Pembengkokan.}
\begin{enumerate}[resume]
  \item Molekul linear yang terletak sepanjang $z$ dibengkokkan pada bidang
        $yz$. \omterm{def:b2:diatomic-mos:bonding}{Orbital nonikatan} manakah yang mulai bertumpang-tindih dengan
        $h_1 + h_2$?
  \item Apa yang terjadi pada tingkat terisi yang dibawanya?
  \item Tingkat ikatan manakah yang naik sedikit, dan mengapa?
  \item Simpulkan: mengapa air bengkok?
  \item Bandingkan sudut terukur dengan $90^\circ$ (ikatan p murni) dan
        $109.5^\circ$.
  \item \omterm{def:b2:diatomic-mos:bonding}{Orbital nonikatan} manakah yang tidak tersentuh oleh pembengkokan?
\end{enumerate}

\medskip
\noindent\textbf{Bagian IV --- Spektrum fotoelektron.}
\begin{enumerate}[resume]
  \item Berapa pita valensi yang ditunjukkan oleh spektrum?
  \item Menurut pendekatan Koopmans, dari orbital mana pita pertama berasal?
  \item Mengapa pita itu sempit, sedangkan pita orbital-orbital ikatan
        lebar?
  \item Bandingkan energi ionisasi pertama air dengan energi ionisasi
        pertama atom oksigen, lalu jelaskan perbedaannya.
  \item Apa yang terjadi dengan ``dua pasangan elektron bebas yang setara''
        pada struktur Lewis?
  \item Nyatakan jumlah pita ionisasi valensi air dan energi pita yang
        terendah.
\end{enumerate}
\end{problem}
